\documentclass[11pt,twoside]{article}
\usepackage{./asp2023}
\aspSuppressVolSlug
\resetcounters

\bibliographystyle{asp2021}

\markboth{G. Mustafa, N. Parmiggiani, C. Sartori, A. Bulgarelli}{Machine-Learning Prediction of the COSI Background}

\urlstyle{same}
\begin{document}

\title{Machine Learning Models for Predicting the COSI Background for Gamma-Ray Burst Detection}

\author{G.~Mustafa,$^{1,2}$ N.~Parmiggiani,$^{2}$ C.~Sartori,$^{1}$ A.~Bulgarelli,$^{2}$ F.~Fenu,$^{3}$ 
V.~Fioretti,$^{2}$ E.~Cavazzuti,$^{3}$ E.~Burns,$^{11}$ L.~Castaldini,$^{2}$ L.~Di~Gesu,$^{3}$
A.~Di~Piano,$^{2}$ R.~Falco,$^{2}$ A.~Femia,$^{2}$ S.~Gallego,$^{7}$ P.~Janowski,$^{7}$ J.~Lommler,$^{7}$
M.~Maniscalco,$^{2,4}$ I.~Martinez-Castellanos,$^{5,6}$ U.~Oberlack,$^{7}$ G.~Panebianco,$^{2}$ A.~Rizzo,$^{12}$
S.~E.~Boggs,$^{8}$ D.~H.~Hartmann,$^{9}$ C.~A.~Kierans,$^{5}$ J.~A.~Tomsick,$^{10}$
and A.~Zoglauer$^{10}$}

\affil{$^{1}$Department of Computer Science and Engineering, University of Bologna, Italy}

\affil{$^{2}$INAF/OAS Bologna, Bologna, Italy}

\affil{$^{3}$Agenzia Spaziale Italiana (ASI), Rome, Italy}

\affil{$^{4}$University of Padova, Padova, Italy}

\affil{$^{5}$NASA Goddard Space Flight Center, Greenbelt, MD, USA}

\affil{$^{6}$Department of Astronomy, University of Maryland, College Park, MD, USA}

\affil{$^{7}$Institute of Physics, Johannes Gutenberg University Mainz, Germany}

\affil{$^{8}$Astronomy and Astrophysics, University of California, San Diego, CA, USA}

\affil{$^{9}$Department of Physics and Astronomy, Clemson University, Clemson, USA}

\affil{$^{10}$Space Sciences Laboratory, University of California, Berkeley, CA, USA}

\affil{$^{11}$Department of Physics and Astronomy, Louisiana State University, LA, USA}

\affil{$^{12}$INAF/OA Catania, Catania, Italy}




\begin{abstract}
The Compton Spectrometer and Imager (COSI) is a NASA Small Explorer mission designed to observe the 0.2--5 MeV gamma-ray sky. Background variations make it difficult to detect transient sources such as gamma-ray bursts (GRBs). We use the COSI Data Challenge 4 (DC4) simulation to study the background for one HEALPix sky direction. Events are selected using the Angular Resolution Measure (ARM) and binned at 1 s to construct a background light curve. We apply supervised machine learning to predict the expected background count rate from spacecraft pointing and orbital information. We compare Poisson regression and Gaussian Mixture Regression (GMR) on a random 30\% evaluation set (1,800,308 1-s bins). Their mean Poisson deviances are 1.0550 and 1.0553, with biases of $-4.9\times10^{-4}$ and $-0.0534$ counts\,s$^{-1}$, respectively. We select Poisson regression because it directly models photon counts and has a lower absolute bias. Its Pearson-residual standard deviation is 1.000 at 1 s. Future work will test the method with simulated GRBs and extend it to blind transient searches across the COSI field of view.
\end{abstract}


\section{Introduction}

The Compton Spectrometer and Imager (COSI) is a wide-field gamma-ray telescope designed to observe the 0.2--5 MeV sky
\citep{Tomsick2023}. We use simulated data from the COSI Germanium Compton Telescope (GCT) to study its background. 
Our goal is to predict the background light curve using supervised machine learning. We use spacecraft pointing and orbital information as inputs and simulated background counts as training targets. We compare Poisson regression and Gaussian Mixture Regression (GMR) to estimate the expected background count rate. These predictions can support future gamma-ray burst (GRB) searches by comparing observed counts with the expected background. We evaluate the models for one HEALPix sky direction as a first step
toward blind GRB searches.

\section{Dataset and Feature Engineering}

We use the three-month COSI Data Challenge~4 (DC4) simulation \citep{Karwin2026,Gallego2026}, produced with MEGAlib \citep{Zoglauer2006}, containing simulated GCT background events and a spacecraft orientation file with time-stamped orbital position and pointing information. HEALPix divides the sky into equal-area pixels \citep{Gorski2005}. Our method is designed to construct a background light curve for each HEALPix direction within the GCT field of view; here, we analyze one reference direction at $N_{\mathrm{side}}=32$. Events are selected using the Angular Resolution Measure (ARM) \citep{Kierans2024}, defined as $\Delta\phi_{\mathrm{ARM}}=\phi^{\mathrm{geo}}-\phi^{\mathrm{kin}}$, where $\phi^{\mathrm{geo}}$ is the geometric scatter angle relative to the tested sky direction and $\phi^{\mathrm{kin}}$ is the Compton scatter angle calculated from the measured energy deposits. For the selected direction, we apply $|\Delta\phi_{\mathrm{ARM}}|\leq13^\circ$ and count the accepted events in 1-s bins to construct its background light curve. 
The $\pm13^\circ$ ARM cut was selected by optimizing the signal-to-noise ratio (SNR) using DC3 Crab source and background simulations. We exclude invalid and SAA-affected intervals with a 7-s guard, together with the first 100 orbits to reduce early activation effects. Event times (\texttt{TimeTags}) are aligned with spacecraft orientation timestamps (\texttt{TimeStamp}) on the 1-s timeline. The model uses 11 spacecraft features: altitude; sine and cosine of the Earth-zenith latitude and longitude; and Cartesian coordinates of the spacecraft X- and Z-axis pointing directions. Five two-orbit blocks are excluded from the random split and reserved for visualization. The remaining 6,001,024 valid 1-s bins are randomly split into training (70\%; 4,200,716) and evaluation (30\%; 1,800,308). Feature standardization is fitted only on the training data. We report model performance on the 30\% evaluation set and use orbits 699--700 for Figure~1. 


\section{Background Models}

We use two machine-learning regression methods implemented in Python with scikit-learn, trained on 4,200,716 one-second background bins. The 11 spacecraft features form the input vector $\mathbf{x}_t$, and $N_t$ is the observed count at time $t$. Poisson regression \citep{McCullagh1989} assumes $N_t\mid\mathbf{x}_t\sim\mathrm{Poisson}(\lambda_t)$, with $\log(\lambda_t)=\beta_0+\boldsymbol{\beta}^{T}\mathbf{x}_t$, where $\lambda_t$ is the expected count and $\beta_0,\boldsymbol{\beta}$ are fitted coefficients. We use \texttt{PoissonRegressor} without regularization.


Four-component Gaussian Mixture Regression (GMR) \citep{Ghahramani1994} models standardized features $\mathbf{x}_t$ and $z_t$, the standardized $\log(1+N_t)$, jointly:
\[
p(\mathbf{x}_t,z_t)=\sum_{k=1}^{4}\pi_k\mathcal{N}\!\left(\begin{bmatrix}\mathbf{x}_t\\z_t\end{bmatrix};\boldsymbol{\mu}_k,\boldsymbol{\Sigma}_k\right).
\]
Here, $\mathcal{N}$ is a multivariate Gaussian density, $k$ indexes the components, $\pi_k$ is the mixture weight, and $\boldsymbol{\mu}_k,\boldsymbol{\Sigma}_k$ are the component mean and covariance. We fit \texttt{GaussianMixture} with full covariance. Conditioning on $\mathbf{x}_t$ gives the log-count mean $m_k$ and variance $v_k$ after reversing standardization. We combine $\exp(m_k+v_k/2)-1$ with conditional mixture weights to obtain $\hat{N}_t=\mathbb{E}[N_t\mid\mathbf{x}_t]$. For both models, $r_t=(N_t-\hat{N}_t)/\sqrt{\hat{N}_t}$ denotes the Pearson residual, where $\hat{N}_t$ is the predicted count ($\lambda_t$ for Poisson regression). 


\section{Results}


We evaluate both models on the random 30\% evaluation set described in Section~2, containing 1,800,308 one-second bins. Poisson regression gives a mean Poisson deviance of 1.0550, a bias of $-4.9\times10^{-4}$ counts\,s$^{-1}$, and a Pearson-residual standard deviation of 1.000. GMR gives 1.0553, $-0.0534$ counts\,s$^{-1}$, and 0.994, respectively (Table~1). We select Poisson regression as the reference model because it directly models photon counts and has a lower absolute bias. Figure~1 shows a representative background light curve from reserved orbits 699--700. SAA intervals divide the light curve into three valid segments. The upper panel shows the illustrative time-dependent $3\sigma$ threshold, $\lambda_t+3\sqrt{\lambda_t}$, computed for each 1-s bin.


\articlefigure[width=0.81\textwidth]{P210_f1.png}{fig:calib}{Background prediction for held-out orbits 699--700. Upper panel: observed 1-s counts (gray), 60-s moving average (blue), Poisson prediction (red), and illustrative $3\sigma$ anomaly threshold (green dashed), $\lambda_t+3\sqrt{\lambda_t}$, computed per 1-s bin. Gray bands mark SAA intervals. Lower panel: Pearson residuals and their 60-s moving average.}



\begin{table}
\caption{Model performance on the random 30\% evaluation set. $N$ is the number of 1-s bins, Bias is the mean observed minus predicted count rate (counts $\mathrm{s}^{-1}$), $\bar{D}$ is mean Poisson deviance, and $\sigma_r$ is Pearson-residual standard deviation.}
\vspace{0.1in}
\centering
\begin{tabular}{l r r r r}
\hline
Model & $N$ & Bias & $\bar{D}$ & $\sigma_r$ \\
\hline
Poisson & 1,800,308 & $-4.9\times10^{-4}$ & 1.0550 & 1.000 \\
GMR & 1,800,308 & $-0.0534$ & 1.0553 & 0.994 \\
\hline
\end{tabular}
\end{table}


\section{Discussion and Conclusions}

On the random 30\% evaluation set, Poisson regression achieves a mean Poisson deviance of 1.0550, a bias of $-4.9\times10^{-4}$ counts\,s$^{-1}$, and a Pearson-residual standard deviation of 1.000. GMR gives a deviance of 1.0553, a bias of $-0.0534$ counts\,s$^{-1}$, and a residual standard deviation of 0.994. Both models show comparable predictive performance, with Pearson-residual standard deviations close to one at 1 s.  We select Poisson regression as the reference model because it directly models photon counts and has a lower absolute bias. Although GMR offers a more flexible probabilistic representation, it shows no clear improvement in the reported metrics. These results support the use of a regression model for predicting the simulated background from spacecraft pointing and orbital information. This study is limited to one HEALPix direction and background-only simulations. The random evaluation split does not establish generalization to independent observing periods, and GRB detection has not yet been tested. Future work will add simulated GRB counts to background light curves and test transient detection by comparing observed counts with model predictions. We will then extend the analysis to other sky directions within the COSI field of view.


\acknowledgments
COSI is a NASA Small Explorer mission supported under NASA contract 80GSFC21C0059. This work is partially supported by the Italian Space Agency under contract ASI/INAF No.~2024-11-HH.0.

\bibliography{P210}

\end{document}